\section{Presentation of the main results} \label{presentation}


The notations for graph polynomials will be the ones in the works of
Brown and collabora\-tors~\cite{Br2009,BrYe09}. To each connected graph~$G$ a matrix $\tilde M_G$ is associated.
The matrix $\tilde M_G$ is a block
matrix with a diagonal block with the Schwinger parameters $x_i$ associated to the edges as diagonal values, two
blocks made of the incidence matrix and minus its transpose and a zero block to complete it. Schematically, it can be written
\begin{equation*}
	\tilde M_G = \begin{pmatrix}
		\operatorname{Diag}(\mathbf{x}) & -{}^t \operatorname{Inc} \\ \operatorname{Inc} & 0
	\end{pmatrix}.
\end{equation*}
This matrix is singular since the sum over the rows indexed by the vertices of the incidence matrix is zero, but the
determinant of any matrix $M_G$ obtained by removing the row and column corresponding to any vertex gives the first
Symanzik polynomial of the graph~$U_G$, independently of the chosen vertex. It is further independent of any of the choices made in the
presentation of the matrix, like the order of the edges and vertices or the orientation of the edges. This polynomial is of degree~$L$, with
$L$ the dimension of the first homology group of the graph, the loop number in the language of physicists.

The objects which will be used to express our results are the polynomials
$U^{I ,J}_K$ defined for~$I$,~$J$ and~$K$ subsets of the edges of the graph
such that $I$ and~$J$ have the same number of elements~$k$. The polynomial
$U^{I ,J}_K$ is the determinant of the matrix deduced from the graph matrix $M_G$ by
removing the rows associated to $I$, the columns associated to $J$ and setting
the variables associated to $K$ to zero and has degree $L-k$. From the symmetry properties of $M_G$,
one deduces that $U^{I ,J}_K = U^{J,I}_K$. These graph polynomials may change sign but are otherwise independent of
the choices made in the presentation of the graph. Since
we will need a definite sign for these expressions, it will be preferable to define
$U^{I ,J}$ as the determinant of the matrix where the rows indexed by $I$ are
replaced with rows with only one non zero value 1 in a position corresponding to a row in
$J$. Expansion along the columns indexed by~$J$ or the rows indexed by~$I$ gives back the previous definition up to
a sign, but now the result only depends on the orientation of the edges indexed by
$I$ and $J $ and the bijection between these subsets
defined by the positions of the 1 in the modified matrix. Schnetz in~\cite{SC2021} proposed to put
an explicit sign in front of the determinant of the reduced matrix. My implicit definition turns out to be fully equivalent to his
explicit one. Both approaches can be useful.

These polynomials, in the special case where $I$ and $J$ have only one element,
give the adjoint elements of the matrix $M_G$, so that the correlation between momenta in rows
$i$ and $j$ is given by $U^{\{i\},\{j\}}/U=U^{i,j}/U$ (we do not keep the braces in the case of
singletons to lighten the notations). The Dodgson identities relate determinants of matrices whose elements
are $U^{i,j}$ for~${i \in I}$ and $j \in J$ to the polynomials $U^{I ,J}$. In Appendix \ref{polyn}, I state these
identities and their proofs, paying due attention to the question of signs.

Our results is most easily described for a vacuum diagram. The case of a propagator diagram can be reduced to this one, as will be seen in next
section.
The graph is replaced by the related completed graph,
obtained by linking the two exterior legs in a simple edge. In this case, the exponent associated to the added edge
must be chosen to make the whole vacuum diagram scale invariant.

We consider a numerator expressed as the product for $s=(s_0,s_1)$ in a set~$S$ of $n$ scalar
products $p_{s_0} \cdot p_{s_1}$. A priori, there are no restrictions to the number of times the momentum
associated to a given edge can appear in this product of scalar products. It should be even possible to have higher
powers of some of the scalar products, but it seems to enter the domain of useless generality.
\begin{Definition}
For a non-empty subset~$S_0$ of~$S$, define $U^{S_0}$ as the mean of the determinants defined using
functions $\veps$ from~$S_0$ to the set of edges of the graph such that
$\veps(s) \in \{s_0,s_1\}$, together with the complementary function
$\tilde\veps$ such that $\{\veps(s),\tilde\veps(s)\} =
\{s_0,s_1\}$
\[
	U^{S_0} = \frac1 {2^{\#(S_0)}} \sum_{\veps} U^{\veps(S_0), \tilde\veps(S_0)}.
\]
\end{Definition}
Since the symmetry of the matrix means changing $\veps$
to $\tilde\veps$ does not change the term, $U^{S_0}$ has at most
$2^{\#(S_0)-1}$ different terms. With these objects, we can easily state our main theorem
\begin{Theorem}\label{fund}
The integrand of a massless vacuum diagram in space-time dimension~$d$ with numerator expressed
as a product of $n$ scalar products can be written as a sum over the set
$\mathcal P_k$ of partitions of~$S$ in $k$~parts, with $k$ taking all values for which this set is
non empty
\begin{gather*}
	I = \sum_{k=1}^{\#(S)} \frac{(-1)^{\#(S)-k}\Gamma\bigl(\frac d2 + k\bigr)}{U^{\frac d2 +k} }
			\sum_{\{S_1,\ldots,S_k\} \in \mathcal P_k} \prod_{j=1}^k U^{S_j}.
\end{gather*}
We can observe that this rational function has homogeneity $-L D/2 - n$.
The value of the residue of the diagram is then the projective integral obtained by multiplying $I$ with the
product~${\prod_i \frac 1{\Gamma(\beta_i)} x_i^{\beta_i-1} {\rm d }x_i^{\vphantom{\beta_i}}}$,
indexed on all edges of $\tilde G$. Since it is a projective integral, any delta function fixing a sum
of~$x_i$ give the same result. The exponents $\beta_i$ associated to the edges must sum up to $L D/2 + n$ to ensure that the integral is
projective.
\end{Theorem}
